\section{境界の楕円modulusをMilnor被覆のPrym部分へ実現する}
\label{sec:boundary-milnor-prym}
前節の楕円曲線は、解消上の四点配置から付随させたものであった。
本節では、その曲線へ実際の有限写像を構成する。
ただし対象は\emph{重み付き主部のMilnor fiberに二次Kummer被覆を施したもの}である。
元のMilnor fiberから同じ底座標を保って写すことには、明示的な平方類の障害がある。
最後に、実際の高重みgermの復元について、整数次数付きdg特異点圏で得られる帰結を区別して述べる。

以下$b=b_n,c=c_n$、$n\ge1$とし、$b$は偶数、$c$は奇数、$\gcd(b,c)=1$を用いる。
$\kappa^2\ne1$、$K=\kappa^2$と置き、
\[
 V_{b,c;\kappa}=\{y^{2b}-2\kappa x^cy^b+x^{2c}=1\},\qquad
 E_\kappa:\ \eta^2=z(z^2-2\kappa z+1)
\]
とする。$C_{b,c;\kappa}$は$V_{b,c;\kappa}$の滑らかな射影模型、
$E_\kappa$では無限遠を原点とする。
重み付き斉次性により$V$は局所Milnor fiberの代数的模型である。
高重み項を持つ実境界の解析的Milnor fiberと無条件に同一視しない。

\subsection{元のfiberでの障害と、最小の二次被覆}
\begin{theorem}[平方類を解消するKummer被覆]
\label{thm:boundary-kummer-prym}
$C=C_{b,c;\kappa}$上の有理関数$z=y^b/x^c$を固定する。
$C\to E_\kappa$でこの$z$を引き戻す有理写像は存在しない。
一方、$\widetilde C$を$\mathbb C(C)(a)$、$a^2=x$の滑らかな射影模型とすると、
\begin{equation}\label{eq:milnor-elliptic-map}
 q:\widetilde C\longrightarrow E_\kappa,\qquad
 z=\frac{y^b}{a^{2c}},\qquad \eta=\frac{y^{b/2}}{a^{3c}}
\end{equation}
は次数$2bc$の有限写像となる。
$\pi:\widetilde C\to C$は連結な二重被覆で、分岐点は$x=0$の$2b$点である。
従って
\[
 g(C)=2bc-b-c,\qquad
 g(\widetilde C)=4bc-b-2c-1,\qquad
 \dim\operatorname{Prym}(\widetilde C/C)=2bc-c-1.
\]
\end{theorem}
\begin{proof}
$Q(z)=z^2-2\kappa z+1$とすると、$V$の方程式は$x^{2c}Q(z)=1$である。
$bp+cq_0=1$を満たす整数$p,q_0$を選び、$t=x^py^{q_0}$と置くと
\[
 x=t^bz^{-q_0},\quad y=t^cz^p,\qquad
 t^{2bc}=\frac{z^{2cq_0}}{Q(z)}.
\]
右辺は$Q$の二根で位数$-1$を持つので、このKummer拡大は連結で次数$2bc$。
四つの底点$0,\infty,\rho_+,\rho_-$上の点数は、それぞれ$2c,2b,1,1$である。
Riemann--Hurwitzから$g(C)=2bc-b-c$を得る。
$\rho_\pm$上の二点では$x$の極位数が$b$であり、$z=\infty$上の$2b$点では
$x$の零位数が$bp+cq_0=1$である。
これは$x=0,y^{2b}=1$という$V$の$2b$点でも直接確認できる。

$E_\kappa$へ持ち上げるための関数の平方は
\begin{equation}\label{eq:milnor-square-obstruction}
 zQ(z)=\frac{y^b}{x^{3c}}
      =x\left(\frac{y^{b/2}}{x^{(3c+1)/2}}\right)^2
\end{equation}
である。$x$は単純零点を持つため平方ではない。これが元の$C$での障害である。
同式により平方根を添加する拡大はちょうど$\mathbb C(C)(\sqrt x)$となり、
表示\eqref{eq:milnor-elliptic-map}を得る。滑らかな射影曲線間の有理写像は延長される。
$\widetilde C\to\mathbb P^1_z$の次数は$4bc$、$E_\kappa\to\mathbb P^1_z$は次数2なので$q$の次数は$2bc$。
$x$の極位数$b$は偶数であるから、$\pi$は無限遠では分岐せず、$2b$個の単純零点だけで分岐する。
再びRiemann--Hurwitzを使い、残る数値を得る。
\end{proof}

\begin{corollary}[楕円曲線のPrym因子]
\label{cor:elliptic-prym-factor}
$E_\kappa$は$\operatorname{Prym}(\widetilde C/C)$のisogeny因子である。
より正確には$q^*$の像はPrymに含まれ、
\[
 e_\kappa=\frac{1}{2bc}q^*q_*
 \in\operatorname{End}^0(J(\widetilde C))
\]
はその像への冪等射である。
従って$H^1(E_\kappa,\mathbb Q)$は$H^1(\widetilde C,\mathbb Q)$の
deck involutionに関する反不変部分のHodge直和因子となる。
\end{corollary}
\begin{proof}
$\delta(a,y)=(-a,y)$は$q$の座標では$(z,\eta)\mapsto(z,-\eta)$となる。
後者は$E_\kappa$の逆元写像なので$\delta^*q^*=-q^*$。
Prymは$\ker(1+\delta^*)^0$であり、連結な$q^*(E_\kappa)$はそこへ入る。
曲線のnormとpullbackについて$q_*q^*=[2bc]$であるから$q^*$は有限核を持ち、
表示した$e_\kappa$は冪等である。
これはPrymの標準的なnorm理論\footnote{
D. Mumford, \emph{Prym Varieties I}, in \emph{Contributions to Analysis},
Academic Press (1974), 325--350, \S\S1--3。
normとpullback、および$P=\ker(1+\delta)^0$については特に印刷頁331を参照。
}のこの被覆への適用であり、整数係数や主偏極を保つ積分解を主張するものではない。
\end{proof}

\subsection{全てのPell指数に共通するCiani四次曲線}
\begin{theorem}[genus三の中間模型と三楕円分解]
\label{thm:universal-ciani-milnor}
$\widetilde C$から、$n$に依存しない滑らかな平面四次曲線
\[
 \mathcal C_\kappa:\quad X^4+Y^4-2\kappa X^2Y^2=Z^4
\]
への次数$bc/2$の写像が、アフィンchartで
\[
 (a,y)\longmapsto(X,Y)=(a^c,y^{b/2})
\]
により与えられる。
$\alpha(X,Y)=(-X,Y)$、$\beta(X,Y)=(X,-Y)$とすると三つの商は楕円曲線であり、
\[
 J(\mathcal C_\kappa)\sim E_i^2\times D_\kappa,\qquad
 D_\kappa:\ T^2=r^4-2\kappa r^2+1,
\]
となる。ここで$E_i$は$j=1728$の楕円曲線、$\sim$は複素数体上のisogenyである。
さらに$\operatorname{Prym}(\widetilde C/C)$は$E_i\times D_\kappa$にisogenousな
二次元の部分アーベル多様体を持つ。
\end{theorem}
\begin{proof}
式への代入で写像を得る。一般の$X,Y\ne0$には、$a^c=X$の$c$通りと
$y^{b/2}=Y$の$b/2$通りの持上げがあり、次数は$bc/2$。
$\kappa^2\ne1$で平面四次は滑らかなので、そのgenusは3である。

$\alpha$商と$\beta$商はそれぞれ
\[
 (X^2-\kappa Y^2)^2=1+(K-1)Y^4,\qquad
 (Y^2-\kappa X^2)^2=1+(K-1)X^4
\]
であり、いずれも複素数体上で$j=1728$となる。
$\alpha\beta$商は$r=Y/X,T=1/X^2$による$D_\kappa$である。
三つの商写像による$H^0(\Omega^1)$の引戻しは、Klein四元群の相異なる三つの非自明指標に属し、
各次元1で全体を張る。従って三楕円曲線からのJacobianのpullbackの和はisogenyとなる。
これはCiani四次の標準的な三楕円分解の特殊例である。\footnote{
E. W. Howe, F. Lepr\'evost, B. Poonen,
\emph{Large torsion subgroups of split Jacobians of curves of genus two or three},
Forum Math. 12 (2000), 315--364,
\href{https://arxiv.org/abs/math/9809210}{arXiv:math/9809210}, \S4.1。
本稿では三つの商の方程式と文字分解を直接確認する。
}

$\widetilde C\to\mathcal C_\kappa$は$\delta$と$\alpha$に関して等変である。
$\alpha$反不変部分は$\beta$商と$\alpha\beta$商の二楕円因子であり、
そのpullbackは有限核を持って$\operatorname{Prym}(\widetilde C/C)$へ入る。
これが最後の主張である。
\end{proof}

$D_\kappa$から前節の$E_\kappa$へは
\begin{equation}\label{eq:ciani-level2-isogeny}
 (r,T)\longmapsto(z,\eta)=(r^2,rT)
\end{equation}
という次数2の写像がある。
そのdeck変換$(r,T)\mapsto(-r,-T)$は固定点を持たず、
原点を適合させれば二同種写像となる。
\[
 j(D_\kappa)=\frac{64(K+3)^3}{(K-1)^2}
 =j\left(E_{\kappa^\vee}\right),\qquad
 (\kappa^\vee)^2=\frac K{K-1}.
\]
従って前節のlevel $2$ isogenyは、ここでは\emph{共通のgenus三模型の商写像}として実現される。
$n=1$では$g(C)=13$、$g(\widetilde C)=27$、Prym次元14、
$\widetilde C\to\mathcal C_\kappa$の次数5、$\widetilde C\to E_\kappa$の次数20である。
高genusの増大と独立に、可変な楕円因子は同じ式で残る。

\subsection{Prym部分の周期はLegendre型方程式を満たす}
\begin{proposition}[可変楕円因子のPicard--Fuchs方程式]
\label{prop:boundary-prym-PF}
$E_\kappa$上の形式$\omega_\kappa=dz/\eta$と平坦に動くcycleの周期$\Pi$は
\begin{equation}\label{eq:boundary-prym-PF}
 (\kappa^2-1)\Pi''+2\kappa\Pi'+\tfrac14\Pi=0
\end{equation}
を満たす。同式は系\ref{cor:elliptic-prym-factor}で得たPrym部分のrank二の周期系にも適用される。
$q=(1-\kappa)/2$とすれば、これは
$q(1-q)\Pi_{qq}+(1-2q)\Pi_q-\Pi/4=0$、すなわち
パラメータ$(1/2,1/2;1)$のGauss超幾何方程式である。
\end{proposition}
\begin{proof}
$P=z(z^2-2\kappa z+1)$と書く。
$\partial_\kappa\omega=z^2dz/P^{3/2}$、
$\partial_\kappa^2\omega=3z^4dz/P^{5/2}$なので、直接計算で
\[
 \left((\kappa^2-1)\partial_\kappa^2+2\kappa\partial_\kappa+\tfrac14\right)
 \omega_\kappa
 =d_z\left(\frac{z^2(1-z^2)}{2P^{3/2}}\right)
\]
を得る。右辺は楕円曲線の有理関数の微分であり、閉cycleでの積分は零となる。
有限写像によるpullbackとtraceはGauss--Manin接続に適合するため、Prymの直和因子にも同じ式を得る。
\end{proof}

これは$\kappa$という\emph{解析modulusを動かす}周期系である。
固定$\kappa$で値$F$を一周する有限位数のMilnor monodromyや、
元のLG三次曲面のlevel $h$を動かしたarcsine周期系とは異なる。
前節で$j$やisogenyが二周期MF模型を完全には分類しなかったこととも矛盾しない。
ここで回収したのは、主部のMilnor被覆の\emph{有理Hodge因子}であって、圏の同値ではない。

\subsection{実germの高重み情報を整数次数付きdg圏で検出する}
最後に実際の$\mu=27,\tau=24$例へ戻る。
$R_f=\mathbb C[[x,y]]/(f)$とし、
\[
 \mathcal S_f=D^b_{\rm dg}(\operatorname{mod}R_f)/\operatorname{Perf}_{\rm dg}(R_f)
\]
を、その通常の$\mathbb Z$-graded dg enhancementを持つ特異点圏とする。
本節のHochschild cohomologyは、この整数次数付き圏の絶対的なものとする。

\begin{corollary}[実境界と主部のdg特異点圏は非同値]
\label{cor:actual-germ-Zdg-distinction}
前節の一反復germ $f$と、その重み付き主部$F$について
\[
 \dim\mathrm{HH}^0(\mathcal S_f)=24,\qquad
 \dim\mathrm{HH}^0(\mathcal S_F)=27.
\]
従って$\mathcal S_f$と$\mathcal S_F$はdg Morita非同値である。
さらに、同じ$n$の許容する実境界germ同士で$\kappa^2\ne(\kappa')^2$なら、
それらの整数次数付きdg特異点圏もMorita非同値である。
\end{corollary}
\begin{proof}
Hua--Kellerの復元定理\footnote{
Z. Hua, B. Keller, \emph{Cluster categories and rational curves},
Geom. Topol. 28 (2024), 2569--2634,
\href{https://doi.org/10.2140/gt.2024.28.2569}{doi:10.2140/gt.2024.28.2569};
\href{https://arxiv.org/abs/1810.00749v6}{arXiv:1810.00749v6}, Theorem 5.9, pp.36--37。
ここで使うのは一般の孤立超曲面についての定理であり、
三次元flopやNCCRの存在を仮定する後続の定理ではない。
}は、孤立超曲面germについて
$\mathrm{HH}^0(\mathcal S_f)\cong\mathbb C[[x,y]]/(f,f_x,f_y)$を代数として与える。
前節の厳密なTjurina数の計算とMorita不変性から最初の主張が従う。
後半では、Morita同値ならTjurina環が同型となり、Mather--Yauによって接触型が一致する。
前節で証明した実二枝の解析的不変量$\kappa^2$が一致しなければならない。
\end{proof}

\begin{remark}[基礎環・次数・実germを保つ]
二周期MFの相対的Hochschild環をJacobian環と同定する公式と、
上の整数次数付きdg特異点圏の絶対的$\mathrm{HH}^0$をTjurina環と同定する公式を入れ替えない。
本系はHua--Kellerの指定するenhancementでの結論であり、
本稿ではperiodization・忘却の全比較を構成せず、一般の二周期MFの復元へ無条件に言い換えない。
同様に、Kummer--Prym構成は重み付き模型のMilnor curveについて証明したもので、
$\tau=24$の実際のgermに解析的に移したものではない。

Prymのnorm理論、Cianiの三楕円分解、超幾何方程式、dg復元の一般定理は既知である。
本稿の位置付けは、境界から得た同じ$K$を明示的な有限写像とHodge因子へ結び、
何を付加すれば幾何学的対応が実現し、どの範囲にはまだ移せないかを示すことにある。
Prym全体の偏極を保つ直積分解、実germのPrym比較、HMSは未構成である。
\end{remark}
